\documentclass{situationist-preprint}

\title{Inscrutability Reconsidered}
\msid{LNG-16-b}
\msfamily{LNG-16}
\msdiscipline{Philosophy of Language}
\mstype{Companion manuscript}
\mskeywords{inscrutability of reference; proxy functions; causal constraint; theoretical terms; Davidson; Lewis}

\begin{document}
\maketitle

\begin{abstract}
LNG-16-a argues that reference is scrutable because a causal constraint on observational applications excludes every proxy reinterpretation that moves the salient causes of those applications. Two challenges remain. A proxy function could leave observable objects fixed and permute only unobservables, so that theoretical terms would be inscrutable; and Davidson held that inscrutability is harmless, a mere relativity to a scheme of translation. We argue that theoretical terms are applied observationally through instruments, so that the causal constraint reaches them, and that Davidson's own later account of interpretation, which gives the shared cause of responses a role in fixing content, contains the constraint and so is incompatible with his earlier thesis.
\end{abstract}

\section{Permutations of unobservables}

Proposition 2.2 of LNG-16-a states that no interpretation that satisfies the causal constraint, Principle P1 of LNG-16-a, remains correct under a permutation that moves the salient causes of observational applications. A permutation that leaves every observable object in place satisfies that hypothesis vacuously. Let $\pi$ exchange each electron with some other unobservable, keeping fixed every object that anyone has ever perceived. Then every sentence keeps its truth value under the reinterpretation by $\pi$, and the causal constraint, applied only to perceived objects, appears silent. Inscrutability seems to return for the theoretical vocabulary.

It does not, because terms such as \enquote{electron} are applied observationally, though not in response to perception of electrons alone. A physicist who looks at a cloud chamber and says \enquote{there goes an electron} applies the term in response to a track, and the salient cause of the track, the object whose passage explains why there is a track there and not elsewhere, is the electron. Principle P1 of LNG-16-a speaks of the object that is the salient cause of the application, not of the object perceived. The causal chain from electron to track to utterance is as determinate as the chain from rabbit to retina to utterance; the instrument lengthens it without making its first link indeterminate. So $\pi$ moves the salient cause of observational applications of \enquote{electron}, and by Proposition 2.2 of LNG-16-a it yields an incorrect interpretation. The same holds of any term with instrument-mediated applications, which in the sciences includes almost every term for an unobservable.

\begin{objection}
The salient cause of a track is a matter of explanatory interest, and what is salient depends on the theory. The constraint therefore presupposes the very theory whose interpretation is in question \citep{Putnam1980}.
\end{objection}

\begin{reply}
Salience is a feature of the causal structure: the electron makes the difference to where the track is, and the electron's mother atom does not. That a theory is needed to know which object made the difference does not make the fact that it did a part of the theory. LNG-16-a replies to the general form of the objection by distinguishing a condition on the reference relation from a sentence to be interpreted; the instrumental case is an instance.
\end{reply}

\section{Davidson's relativity}

\citet{Davidson1979} accepted the permutation argument but drew a milder conclusion: reference is relative to a scheme of interpretation, and since every scheme that preserves truth conditions is equally good, nothing is lost. His later work changed the evidence base. On the triangulation account, the content of a speaker's observational sentences is fixed by the common cause of her responses and her interpreter's responses to it \citep{Davidson2001}. That is a version of the causal constraint, and it violates Assumption 2.1 of LNG-16-a, which Davidson's earlier thesis required: the facts that fix reference are no longer exhausted by truth conditions and dispositions to assent. A proxy reinterpretation that sends \enquote{rabbit} to the cosmic complement of each rabbit leaves the truth conditions intact but does not leave the common cause intact; the interpreter and the speaker were both responding to the rabbit. So by Davidson's later lights the reinterpretation is not equally good, and relativity to a scheme reduces to the choice of notation.

\citet{Lewis1984} reached scrutability by a different route, appealing to the greater eligibility of natural properties as referents. The causal constraint needs no such ranking: it excludes proxy reinterpretations whether or not their referents are natural. The two routes are compatible, and the second is the less demanding.

\section{Conclusion}

The causal constraint covers theoretical terms through their instrumental applications and is present in Davidson's own mature theory of interpretation. Inscrutability does not survive for the vocabulary of science, and its harmless form survives only as a choice among notations.

\begin{thebibliography}{9}\small
\bibitem[Davidson(1979)]{Davidson1979} Davidson, D. (1979). The inscrutability of reference. \emph{The Southwestern Journal of Philosophy}, 10(2), 7--19.
\bibitem[Davidson(2001)]{Davidson2001} Davidson, D. (2001). \emph{Subjective, Intersubjective, Objective}. Clarendon Press.
\bibitem[Lewis(1984)]{Lewis1984} Lewis, D. (1984). Putnam's paradox. \emph{Australasian Journal of Philosophy}, 62(3), 221--236.
\bibitem[Putnam(1980)]{Putnam1980} Putnam, H. (1980). Models and reality. \emph{The Journal of Symbolic Logic}, 45(3), 464--482.
\end{thebibliography}

\end{document}
